\documentclass[12pt, a4paper]{article}
 % article, report, book, beamer
% tikz
\usepackage{tikz}
\usetikzlibrary{positioning, arrows.meta, shapes}
% colors
\usepackage{xcolor}
\usepackage{titlesec}
% header
\usepackage{fancyhdr}
\usepackage[utf8]{inputenc}
% Basic
\usepackage{amsfonts}
\usepackage{amssymb}
\usepackage{amsmath}

\title{Exercise 02}
\author{LoveKapibarasan}
\date{\today}                   % or specific date

\begin{document}
\maketitle                      % Create title
% Content here
\pagestyle{fancy}
\fancyhf{} 
\fancyhead[L]{\leftmark}     % section name
\fancyhead[R]{\thepage}      % page number
\fancyhead[C]{Answer}   % Custom number
% colors
\titleformat{\section}
  {\color{red}\normalfont\Large\bfseries}  % Blue, Bold
  {\thesection}{1em}{}
\titleformat{\subsection}
  {\color{red}\normalfont\Large\bfseries}  % Blue, Bold
  {\thesection}{1em}{}
%=====

\section* {2. Finite and infinite probability space}

\subsection* {E1}
\begin{enumerate}
  \item $ \underline{\underline{ \frac{1}{2} }} $
  \item $1 - \left( \frac{1}{2} \right) ^ 3 = \underline{\underline{ \frac{7}{8}  }}$
  \item $ \frac{4}{12} = \underline{\underline{ \frac{1}{3} }}$
\end{enumerate}

\subsection* {E2}
\begin{enumerate}
  \item $\underline{\underline{ \frac{7}{15} }} $
  \item $\underline{\underline{ \frac{4}{15}  }}$
  \item $\frac{ | \{ 1, 2, 3, 4, 6, 8, 10. 12, 14 \} | }{12} = \underline{\underline{ \frac{9}{12} }}$
  \item $\frac{ | \{  2,  4 \} | }{15} = \underline{\underline{ \frac{2}{15} }}$
\end{enumerate}

\subsection* {E3}
\begin{enumerate}
  \item $\frac{1 \cdot 4}{ C(5, 2) } = \underline{\underline{ \frac{2}{5} }} $
\end{enumerate}

\subsection* {E4}
\begin{align*}
|man| &= 10 \\
|man \cap brown| &= 5 \\
|brown| &= 30 / 2 = 15
\end{align*}
\begin{enumerate}
  \item $\frac{10 + 15 - 5}{ 30 }= \underline{\underline{ \frac{2}{3} }} $
\end{enumerate}

\subsection* {E5}
$$
P(T)+P(H) = 1 \\
P(H) = 2P(T)
$$
\begin{itemize}
  \item $ P(H) = \underline{\underline{ \frac{1}{3} }} $
  \item $ P(T) = \underline{\underline{ \frac{2}{3} }} $
\end{itemize}


\subsection* {E6}
\begin{enumerate}
  \item $1 - P(A) = \underline{\underline{ 0.4 }} $
  \item $\underline{\underline{ 0.7 }} $
  \item $P(A) + P(B) - P (A \cap B)=\underline{\underline{ 0.7 }} $
  \item $\underline{\underline{ 0.3 }} $
\end{enumerate}

\subsection* {E7}
\begin{enumerate}
  \item $p_2 + p_4 + p_6 = \underline{\underline{ 0.6 }} $
  \item $\underline{\underline{ 0.7 }} $
  \item $p_1 + p_2 = \underline{\underline{ 0.4 }} $
  \item $\underline{\underline{ 0 }} $
\end{enumerate}

\subsection* {E8}
\begin{enumerate}
  \item $1 - \left( P(a_2) + P(a_3) + P(a_4) \right) = \underline{\underline{ 0.3 }} $
  \item
	\begin{align*}
	P(a_1) + P(a_2) + P(a_3) + P(a_4) &= 1 \\
	P(a_1) &= 3P(a_2) \\
	P(a_3) = P(a_4) &= 0.2 \\
	P(a_2) &= \underline{\underline{\frac{3}{20}}} \\
	P(a_1) &= \underline{\underline{\frac{9}{20}}}
	\end{align*}
\end{enumerate}


\subsection* {E9}
For all $y = l$, $x \in [k, k + 1)$ ($l \in \mathbb{R}$, $k \in \mathbb{Z}$), the center of coin $x_0$ must satisfy $\frac{k}{4} < x_0 < \frac{3k}{4}$ in a length 1 line.

$$\underline{\underline{\frac{1}{2}}}$$


%=====
\end{document}
